\subsubsection*{Solution}
\paragraph*{Step-by-Step Derivation}

Let the normalized Choi state be denoted by $\gamma_{q,d}$, with all logarithms taken to base $2$, so that capacities are measured in qubits per channel use. We assume $d\geq 2$.

\subparagraph*{1. Simplification at the special value of $q$}

The symmetric and antisymmetric subspace dimensions are

\[
d_{\mathrm{sym}}=\frac{d(d+1)}{2}, \qquad d_{\mathrm{asym}}=\frac{d(d-1)}{2}.
\]

For

\[
q=\frac{d+1}{2d}, \qquad 1-q=\frac{d-1}{2d},
\]

the coefficients simplify to

\[
\frac{q}{d_{\mathrm{sym}}} = \frac{(d+1)/(2d)}{d(d+1)/2} = \frac{1}{d^2},
\]

and similarly,

\[
\frac{1-q}{d_{\mathrm{asym}}} = \frac{(d-1)/(2d)}{d(d-1)/2} = \frac{1}{d^2}.
\]

Therefore,

\[
\gamma = \frac{1}{d^2} \left( \psi_+\otimes P_{\mathrm{sym}} + \psi_-\otimes P_{\mathrm{asym}} \right),
\]

where $\psi_\pm=|\psi_\pm\rangle\langle\psi_\pm|$.

Using

\[
P_{\mathrm{sym}}=\frac{I+F}{2}, \qquad P_{\mathrm{asym}}=\frac{I-F}{2},
\]

with $F$ the swap operator, and

\[
\psi_++\psi_- = |00\rangle\langle00|+|11\rangle\langle11|,
\]

\[
\psi_+-\psi_- = |00\rangle\langle11|+|11\rangle\langle00|,
\]

we obtain

\[
\gamma = \frac{1}{2d^2} \left[ \left( |00\rangle\langle00|+|11\rangle\langle11| \right)\otimes I + \left( |00\rangle\langle11|+|11\rangle\langle00| \right)\otimes F \right].
\]

This is the special private channel analyzed in Sec. IV-A of \href{\detokenize{https://arxiv.org/html/2510.04527v2}}{Wu and Wang}.

\subparagraph*{2. Explicit form of the channel}

Write an arbitrary input state on $\mathbb C^2\otimes\mathbb C^d$ in block form:

\[
\rho=
\begin{pmatrix}
X_{00}&X_{01}\\
X_{10}&X_{11}
\end{pmatrix},
\qquad
X_{rs}\in\mathcal B(\mathbb C^d).
\]

Applying the normalized Choi reconstruction formula gives

\[
{
\mathcal N(\rho)
=
\frac{1}{d}
\begin{pmatrix}
\operatorname{Tr}(X_{00})I_d&X_{01}^{T}\\
X_{10}^{T}&\operatorname{Tr}(X_{11})I_d
\end{pmatrix}.
}
\]

One Kraus representation is

\[
K_{ij} = \frac{1}{\sqrt d} \left( |0j\rangle\langle0i| + |1i\rangle\langle1j| \right), \qquad i,j=0,\ldots,d-1.
\]

Indeed,

\[
\sum_{i,j}K_{ij}^{\dagger}K_{ij} = I_{2d}.
\]

The associated complementary channel is

\[
{ \mathcal N^c(\rho) = \frac{1}{d} \left( X_{00}\otimes I_d + I_d\otimes X_{11} \right). }
\]

In particular, the complementary output is independent of the off-diagonal blocks $X_{01}$ and $X_{10}$.

\subparagraph*{3. One-shot coherent information}

The one-shot coherent information is

\[
Q^{(1)}(\mathcal N) = \max_{\rho} \left[ S(\mathcal N(\rho))-S(\mathcal N^c(\rho)) \right].
\]

Since $\mathcal N^c(\rho)$ does not depend on $X_{01}$ and $X_{10}$, while deleting these blocks is a pinching operation that cannot decrease the entropy of $\mathcal N(\rho)$, the optimum can be taken with

\[
X_{01}=X_{10}=0.
\]

Let

\[
p=\operatorname{Tr}X_{00}, \qquad 1-p=\operatorname{Tr}X_{11}.
\]

The channel output is then

\[
\mathcal N(\rho)
=
\frac{1}{d}
\begin{pmatrix}
pI_d&0\\
0&(1-p)I_d
\end{pmatrix}.
\]

Its eigenvalues are $p/d$ and $(1-p)/d$, each with multiplicity $d$. Hence

\[
S(\mathcal N(\rho)) = \log_2 d+h_2(p),
\]

where

\[
h_2(p)=-p\log_2p-(1-p)\log_2(1-p).
\]

For fixed $p$, the complementary entropy is minimized by taking the diagonal blocks to be rank one:

\[
X_{00}=p|\phi_0\rangle\langle\phi_0|, \qquad X_{11}=(1-p)|\phi_1\rangle\langle\phi_1|.
\]

The nonzero eigenvalues of the complementary output are then

\[
\frac{1}{d} \quad\text{once},
\]

\[
\frac{p}{d} \quad\text{with multiplicity }d-1,
\]

and

\[
\frac{1-p}{d} \quad\text{with multiplicity }d-1.
\]

Consequently,

\[
S(\mathcal N^c(\rho)) = \log_2d+\frac{d-1}{d}h_2(p).
\]

The coherent information becomes

\[
I_c(\rho,\mathcal N) = \left(\log_2d+h_2(p)\right) - \left(\log_2d+\frac{d-1}{d}h_2(p)\right) = \frac{h_2(p)}{d}.
\]

Since

\[
\max_{0\leq p\leq1}h_2(p)=1
\]

at $p=1/2$, we find the unconditional one-shot result

\[
{ Q^{(1)}(\mathcal N)=\frac{1}{d}. }
\]

\subparagraph*{4. Regularized quantum capacity}

The quantum capacity is the regularized coherent information:

\[
Q(\mathcal N) = \lim_{n\rightarrow\infty} \frac{1}{n} Q^{(1)}(\mathcal N^{\otimes n}).
\]

Thus the one-shot calculation alone proves only

\[
Q(\mathcal N)\geq \frac{1}{d}.
\]

The cited private-channel result invokes the Spin Alignment Conjecture to show that, for every $n$,

\[
Q^{(1)}(\mathcal N^{\otimes n}) = \frac{n}{d}.
\]

The conjecture identifies the optimizer with the tensor power of a degradable two-dimensional restriction of the channel. Additivity of coherent information for that degradable restriction then gives

\[
Q(\mathcal N) = \lim_{n\to\infty}\frac{1}{n}\frac{n}{d} = \frac{1}{d}.
\]

The conjectural qualification is important: the full Spin Alignment Conjecture has not been proved in general. A rigorous result currently establishes it only at the two-letter level (\href{\detokenize{https://link.springer.com/article/10.1007/s00023-025-01592-w}}{Alhejji, Annales Henri Poincaré}); moreover, a stronger majorization version has a three-qubit counterexample, although that does not itself disprove the original entropy conjecture (\href{\detokenize{https://arxiv.org/abs/2603.25410}}{Song and Chen}).

Without assuming the conjecture, the transposition bound quoted in the private-channel paper gives

\[
Q(\mathcal N)\leq \log_2\left(1+\frac1d\right),
\]

so the presently rigorous unconditional statement is

\[
\frac1d \leq Q(\mathcal N) \leq \log_2\left(1+\frac1d\right).
\]
